

\subsection{Audit Complexity}











Let $A$ be an auditable algorithm that solves (or claims to solve) DAC
with a known constant-factor approximation, say $k$.  Let
$\mathcal{K}$ denote the certificate generated by $A$ for an input
formula $F$ over $n$ variables.  We expect $\mathcal{K}$ to be a
string of Boolean values of size that is a function of $n$.  We can
now define the notion of an auditor for $A$ as follows.

\begin{definition} An auditor for $A$ is an algorithm $B$ that
  takes as inputs the formula $F$,
  the estimate $\Cest$ returned by $A$, and the
  certificate $\mathcal{K}$.  Algorithm $B$ uses an
  oracle of complexity class $\mathcal{C}$ to certify that $\Cest$
  lies within the claimed $k$-factor approximation of $\sol{F}$.  The {\em
    audit complexity of algorithm $A$} is defined to be the tuple ($\mathcal{C}, r, t$), where the auditor can certify the answer
  computed by $A$ by making $t$ calls to an oracle of complexity class
  $\mathcal{C}$, where each call/query is over a formula with at most
  $r$ variables.
\end{definition}
In our context, motivated by Stockmeyer's result~\cite{stockmeyer}, we restrict
$\mathcal{C}$ to be $\Sigmatwop$, and further fix the number of calls $t$ to be $1$. Hence, we omit $\mathcal{C}$ and $t$, and simply refer to $r$ as the audit complexity of the algorithm $A$.




